% Partie : théorie de la démonstration
% Chapitre : normalisation
% Section : conversions par permutation

\documentclass[../../../include/open-logic-section]{subfiles}

\begin{document}

\olfileid[fr]{pt}{nor}{per}

\olsection{Conversions par permutation}

Les coupures de \Log{N1i} sont des segments qui se
terminent par la prémisse majeure d'une règle
d'élimination. Une étape de la démonstration de
normalisation consiste à raccourcir les coupures
maximales, afin de diminuer la longueur de coupure
de la !!{proof}. Une coupure de longueur~$>1$
peut se terminer de différentes façons :
sa dernière occurrence~$!A$ provient de \Elim{\lor}
ou de $\Elim{\lexists}$, puis sert de prémisse
majeure à une autre élimination.

Par exemple, si les deux règles sont \Elim{\lor}
et \Elim{\land}, on a $!A \ident !B \land !C$,
et une sous-preuve~$\delta_1$ se termine ainsi :
\[
\AxiomC{}
\RightLabel{$\delta_2$}
\DeduceC{$!D \lor !E$}
\AxiomC{$\Discharge{!D}{x}$}
\RightLabel{$\delta_3$}
\DeduceC{$!B \land !C$}
\AxiomC{$\Discharge{!E}{y}$}
\RightLabel{$\delta_4$}
\DeduceC{$!B \land !C$}
\DischargeRule{\Elim{\lor}}{x\,y}
\TrinaryInfC{$!B \land !C$}
\RightLabel{\Elim{\land}[1]}
\UnaryInfC{$!B$}
\DisplayProof
\]
La dernière occurrence~$!A_n$ de la coupure est
la conclusion de \Elim{\lor}, c'est-à-dire
l'occurrence inférieure de~$!B \land !C$.
L'avant-dernière occurrence~$!A_{n-1}$ est une
des prémisses mineures de cette règle.

Pour raccourcir la coupure, nous pouvons
\emph{permuter} les deux éliminations :
nous remplaçons la sous-preuve~$\delta_1$
de~$\delta$ par
\[
\AxiomC{}
\RightLabel{$\delta_2$}
\DeduceC{$!D \lor !E$}
\AxiomC{$\Discharge{!D}{x}$}
\RightLabel{$\delta_3$}
\DeduceC{$!B \land !C$}
\RightLabel{\Elim{\land}[1]}
\UnaryInfC{$!B$}
\AxiomC{$\Discharge{!E}{y}$}
\RightLabel{$\delta_4$}
\DeduceC{$!B \land !C$}
\RightLabel{\Elim{\land}[1]}
\UnaryInfC{$!B$}
\DischargeRule{\Elim{\lor}}{x\,y}
\TrinaryInfC{$!B$}
\DisplayProof
\]
Les coupures qui se terminaient à la prémisse
de \Elim{\land} au-dessous de \Elim{\lor}
se terminent désormais à la prémisse d'une
des deux projections situées au-dessus de
ses prémisses mineures. Chacune perd une
occurrence : sa longueur diminue de~$1$.

Dans cet exemple de projection, les coupures
entièrement contenues dans $\delta_2$,
$\delta_3$ ou~$\delta_4$, ou entièrement
extérieures à~$\delta_1$, conservent leurs
formules, leur rang et leur longueur.
Les éventuelles nouvelles coupures portent
sur la formule projetée, de rang inférieur
à celui de la conjonction. Si la coupure
choisie est maximale, le rang de coupure
ne croît donc pas ; s'il reste le même,
la longueur de coupure diminue.
Ce raisonnement local ne dispense pas de
contrôler les duplications dans les autres
permutations, examinées plus loin.

\begin{prob}
Une permutation raccourcit d'au moins~$1$
un des segments visés. Lorsque \Elim{\land}
est permutée avec \Elim{\lor}, deux segments
de coupure perdent une occurrence ; si leur
rang est maximal et reste le rang de coupure
de la preuve, la longueur de coupure diminue
donc d'au moins~$2$. Une seule permutation
de ce type peut-elle la diminuer de plus
de~$2$ ? Peut-elle faire baisser le
\emph{rang} de coupure ?
\end{prob}

Déplacer une élimination au-dessus d'une
application de \Elim{\lor} ou de
\Elim{\lexists} est appelé une
\emph{conversion par permutation}.
Quelques schémas figurent dans le
\olref{tab:permutation-conversions}.
Les autres sont laissés en exercice.

% \begin{table}
% \begin{multline*}
% \shoveleft{\AxiomC{$!D \lor !E$}
% \AxiomC{$\Discharge{!D}{x}$}
% \shortDeduce
% \DeduceC{$!B \land !C$}
% \AxiomC{$\Discharge{!E}{y}$}
% \shortDeduce
% \DeduceC{$!B \land !C$}
% \DischargeRule{\Elim{\lor}}{x\,y}
% \TrinaryInfC{$!B \land !C$}
% \RightLabel{\Elim{\land}[1]}
% \UnaryInfC{$!B$}
% \DisplayProof}\\
% \shoveright{\longrightarrow\ 
% \AxiomC{$!D \lor !E$}
% \AxiomC{$\Discharge{!D}{x}$}
% \shortDeduce
% \DeduceC{$!B \land !C$}
% \RightLabel{\Elim{\land}[1]}
% \UnaryInfC{$!B$}
% \AxiomC{$\Discharge{!E}{y}$}
% \shortDeduce
% \DeduceC{$!B \land !C$}
% \RightLabel{\Elim{\land}[1]}
% \UnaryInfC{$!B$}
% \DischargeRule{\Elim{\lor}}{x\,y}
% \TrinaryInfC{$!B$}
% \DisplayProof}
% \\
% \shoveleft{\AxiomC{$!D \lor !E$}
% \AxiomC{$\Discharge{!D}{x}$}
% \shortDeduce
% \DeduceC{$!B \lif !C$}
% \AxiomC{$\Discharge{!E}{y}$}
% \shortDeduce
% \DeduceC{$!B \lif !C$}
% \DischargeRule{\Elim{\lor}}{x\,y}
% \TrinaryInfC{$!B \lif !C$}
% \AxiomC{$!B$}
% \RightLabel{\Elim{\lif}}
% \BinaryInfC{$!C$}
% \DisplayProof}\\
% \shoveright{\longrightarrow\ 
% \AxiomC{$!D \lor !E$}
% \AxiomC{$\Discharge{!D}{x}$}
% \shortDeduce
% \DeduceC{$!B \lif !C$}
% \AxiomC{$!B$}
% \RightLabel{\Elim{\lif}}
% \BinaryInfC{$!C$}
% \AxiomC{$\Discharge{!E}{y}$}
% \shortDeduce
% \DeduceC{$!B \lif !C$}
% \AxiomC{$!B$}
% \RightLabel{\Elim{\lif}}
% \BinaryInfC{$!C$}
% \DischargeRule{\Elim{\lor}}{x\,y}
% \TrinaryInfC{$!C$}
% \DisplayProof}
% \\
% \shoveleft{\AxiomC{$!D \lor !E$}
% \AxiomC{$\Discharge{!D}{x}$}
% \shortDeduce
% \DeduceC{$!B \lif !C$}
% \AxiomC{$\Discharge{!E}{y}$}
% \shortDeduce
% \DeduceC{$!B \lif !C$}
% \DischargeRule{\Elim{\lor}}{x\,y}
% \TrinaryInfC{$!B \lif !C$}
% \AxiomC{$!B$}
% \RightLabel{\Elim{\lif}}
% \BinaryInfC{$!C$}
% \DisplayProof}\\
% \shoveright{\longrightarrow\ 
% \AxiomC{$!D \lor !E$}
% \AxiomC{$\Discharge{!D}{x}$}
% \shortDeduce
% \DeduceC{$!B \lif !C$}
% \AxiomC{$!B$}
% \RightLabel{\Elim{\lif}}
% \BinaryInfC{$!C$}
% \AxiomC{$\Discharge{!E}{y}$}
% \shortDeduce
% \DeduceC{$!B \lif !C$}
% \AxiomC{$!B$}
% \RightLabel{\Elim{\lif}}
% \BinaryInfC{$!C$}
% \DischargeRule{\Elim{\lor}}{x\,y}
% \TrinaryInfC{$!C$}
% \DisplayProof}
% \\
% \shoveleft{\AxiomC{$!D \lor !E$}
% \AxiomC{$\Discharge{!D}{x}$}
% \shortDeduce
% \DeduceC{$\lforall[x][!B(x)]$}
% \AxiomC{$\Discharge{!E}{y}$}
% \shortDeduce
% \DeduceC{$\lforall[x][!B(x)]$}
% \DischargeRule{\Elim{\lor}}{x\,y}
% \TrinaryInfC{$\lforall[x][!B(x)]$}
% \RightLabel{\Elim{\lforall}}
% \UnaryInfC{$!B(t)$}
% \DisplayProof}\\
% \shoveright{\longrightarrow\ 
% \AxiomC{$!D \lor !E$}
% \AxiomC{$\Discharge{!D}{x}$}
% \shortDeduce
% \DeduceC{$\lforall[x][!B(x)]$}
% \RightLabel{\Elim{\lforall}}
% \UnaryInfC{$!B(t)$}
% \AxiomC{$\Discharge{!E}{y}$}
% \shortDeduce
% \DeduceC{$\lforall[x][!B(x)]$}
% \RightLabel{\Elim{\lforall}}
% \UnaryInfC{$!B(t)$}
% \DischargeRule{\Elim{\lor}}{x\,y}
% \TrinaryInfC{$!B(t)$}
% \DisplayProof}
% \end{multline*}
% \caption{Quelques conversions par permutation pour \Log{N1i}.}
% \ollabel{tab:permutation-conversions}
% \end{table}

{\small
\begin{longtable}{@{}l@{}}
\AxiomC{$!D \lor !E$}
\AxiomC{$\Discharge{!D}{x}$}
\shortDeduce
\DeduceC{$!B \land !C$}
\AxiomC{$\Discharge{!E}{y}$}
\shortDeduce
\DeduceC{$!B \land !C$}
\DischargeRule{\Elim{\lor}}{x\,y}
\TrinaryInfC{$!B \land !C$}
\RightLabel{\Elim{\land}[1]}
\UnaryInfC{$!B$}
\DisplayProof
$\longrightarrow$\\[6ex]
\quad\AxiomC{$!D \lor !E$}
\AxiomC{$\Discharge{!D}{x}$}
\shortDeduce
\DeduceC{$!B \land !C$}
\RightLabel{\Elim{\land}[1]}
\UnaryInfC{$!B$}
\AxiomC{$\Discharge{!E}{y}$}
\shortDeduce
\DeduceC{$!B \land !C$}
\RightLabel{\Elim{\land}[1]}
\UnaryInfC{$!B$}
\DischargeRule{\Elim{\lor}}{x\,y}
\TrinaryInfC{$!B$}
\DisplayProof
\\[10ex]
\AxiomC{$!D \lor !E$}
\AxiomC{$\Discharge{!D}{x}$}
\shortDeduce
\DeduceC{$!B \lif !C$}
\AxiomC{$\Discharge{!E}{y}$}
\shortDeduce
\DeduceC{$!B \lif !C$}
\DischargeRule{\Elim{\lor}}{x\,y}
\TrinaryInfC{$!B \lif !C$}
\AxiomC{$!B$}
\RightLabel{\Elim{\lif}}
\BinaryInfC{$!C$}
\DisplayProof
$\longrightarrow$\\[6ex]
\AxiomC{$!D \lor !E$}
\AxiomC{$\Discharge{!D}{x}$}
\shortDeduce
\DeduceC{$!B \lif !C$}
\AxiomC{$!B$}
\RightLabel{\Elim{\lif}}
\BinaryInfC{$!C$}
\AxiomC{$\Discharge{!E}{y}$}
\shortDeduce
\DeduceC{$!B \lif !C$}
\AxiomC{$!B$}
\RightLabel{\Elim{\lif}}
\BinaryInfC{$!C$}
\DischargeRule{\Elim{\lor}}{x\,y}
\TrinaryInfC{$!C$}
\DisplayProof
\\[10ex]
\AxiomC{$!D \lor !E$}
\AxiomC{$\Discharge{!D}{x}$}
\shortDeduce
\DeduceC{$!B \lif !C$}
\AxiomC{$\Discharge{!E}{y}$}
\shortDeduce
\DeduceC{$!B \lif !C$}
\DischargeRule{\Elim{\lor}}{x\,y}
\TrinaryInfC{$!B \lif !C$}
\AxiomC{$!B$}
\RightLabel{\Elim{\lif}}
\BinaryInfC{$!C$}
\DisplayProof
$\longrightarrow$\\[6ex]
\AxiomC{$!D \lor !E$}
\AxiomC{$\Discharge{!D}{x}$}
\shortDeduce
\DeduceC{$!B \lif !C$}
\AxiomC{$!B$}
\RightLabel{\Elim{\lif}}
\BinaryInfC{$!C$}
\AxiomC{$\Discharge{!E}{y}$}
\shortDeduce
\DeduceC{$!B \lif !C$}
\AxiomC{$!B$}
\RightLabel{\Elim{\lif}}
\BinaryInfC{$!C$}
\DischargeRule{\Elim{\lor}}{x\,y}
\TrinaryInfC{$!C$}
\DisplayProof
\\[10ex]
\AxiomC{$!D \lor !E$}
\AxiomC{$\Discharge{!D}{x}$}
\shortDeduce
\DeduceC{$\lforall[x][!B(x)]$}
\AxiomC{$\Discharge{!E}{y}$}
\shortDeduce
\DeduceC{$\lforall[x][!B(x)]$}
\DischargeRule{\Elim{\lor}}{x\,y}
\TrinaryInfC{$\lforall[x][!B(x)]$}
\RightLabel{\Elim{\lforall}}
\UnaryInfC{$!B(t)$}
\DisplayProof
\quad$\longrightarrow$\\[8ex]
\quad\AxiomC{$!D \lor !E$}
\AxiomC{$\Discharge{!D}{x}$}
\shortDeduce
\DeduceC{$\lforall[x][!B(x)]$}
\RightLabel{\Elim{\lforall}}
\UnaryInfC{$!B(t)$}
\AxiomC{$\Discharge{!E}{y}$}
\shortDeduce
\DeduceC{$\lforall[x][!B(x)]$}
\RightLabel{\Elim{\lforall}}
\UnaryInfC{$!B(t)$}
\DischargeRule{\Elim{\lor}}{x\,y}
\TrinaryInfC{$!B(t)$}
\DisplayProof
\\
\caption{Quelques conversions par permutation pour \Log{N1i}.}
\ollabel{tab:permutation-conversions}
\end{longtable}
}


\begin{prob}
Donnez les conversions par permutation
pour \Elim{\lor} suivie de \Elim{\lor}
ou de \Elim{\lnot}.
\end{prob}

\begin{prob}
Donnez la conversion par permutation
pour \Elim{\lexists} suivie de
\Elim{\lexists}. Expliquez pourquoi
elle respecte les conditions sur les
constantes propres.
\end{prob}

Il faut également vérifier qu'une permutation
respecte les conditions de fraîcheur.
Considérons
\[
\AxiomC{}
\RightLabel{$\delta_2$}
\DeduceC{$!D \lor !E$}
\AxiomC{$\Discharge{!D}{x}$}
\RightLabel{$\delta_3$}
\DeduceC{$\lexists[x][!B(x)]$}
\AxiomC{$\Discharge{!E}{y}$}
\RightLabel{$\delta_4$}
\DeduceC{$\lexists[x][!B(x)]$}
\DischargeRule{\Elim{\lor}}{x\,y}
\TrinaryInfC{$\lexists[x][!B(x)]$}
\AxiomC{$\Discharge{!B(c)}{z}$}
\RightLabel{$\delta_5$}
\DeduceC{$!C$}
\DischargeRule{\Elim{\lexists}}{z}
\BinaryInfC{$!C$}
\DisplayProof
\]
Si nous remplaçons cet arbre par
\[
\AxiomC{}
\RightLabel{$\delta_2$}
\DeduceC{$!D \lor !E$}
\AxiomC{$\Discharge{!D}{x}$}
\RightLabel{$\delta_3$}
\DeduceC{$\lexists[x][!B(x)]$}
\AxiomC{$\Discharge{!B(c)}{z}$}
\RightLabel{$\delta_5$}
\DeduceC{$!C$}
\DischargeRule{\Elim{\lexists}}{z}
\BinaryInfC{$!C$}
\AxiomC{$\Discharge{!E}{y}$}
\RightLabel{$\delta_4$}
\DeduceC{$\lexists[x][!B(x)]$}
\AxiomC{$\Discharge{!B(c)}{z}$}
\RightLabel{$\delta_5$}
\DeduceC{$!C$}
\DischargeRule{\Elim{\lexists}}{z}
\BinaryInfC{$!C$}
\DischargeRule{\Elim{\lor}}{x\,y}
\TrinaryInfC{$!C$}
\DisplayProof
\]
la constante propre~$c$ pourrait sembler
apparaître dans~$!D$. Dans le premier
arbre, cette hypothèse est déchargée avant
l'élimination existentielle ; dans le
second, elle ne l'est qu'après les deux
applications de cette règle. Si la
constante y apparaissait, la nouvelle
preuve serait incorrecte.

Nous travaillons toutefois avec des preuves
\emph{régulières}. Une constante propre
n'est propre qu'à une seule inférence
et n'apparaît que dans la sous-preuve
correspondante : au-dessus de la prémisse
de \Intro{\lforall}, ou au-dessus de
la prémisse mineure de \Elim{\lexists}.
Ici, $c$ ne peut donc apparaître ni
dans~$\delta_3$ ni dans~$\delta_4$,
et en particulier pas dans les hypothèses
temporaires des deux branches.
Le dessin réutilise le même nom dans les
deux copies. Pour rendre la preuve obtenue
régulière, on renomme séparément les
constantes propres de chaque copie, y
compris celles d'éventuelles inférences
intérieures. On renomme aussi les classes
déchargées lorsque c'est nécessaire pour
éviter toute capture d'une hypothèse
ouverte par une autre règle.

Cet exemple montre enfin une difficulté
pour la mesure de la preuve. La nouvelle
preuve contient deux copies de~$\delta_5$.
Si celle-ci contient une coupure maximale,
sa duplication peut empêcher la diminution
de la longueur de coupure. Le choix de
la coupure doit donc garantir qu'aucune
coupure maximale n'est dupliquée.

La stratégie de normalisation emploiera une
coupure maximale \emph{supérieure et située
le plus à droite}, avec les conditions
précisées dans le théorème. Elles excluent
les coupures maximales dans les sous-preuves
dupliquées et les segments maximaux passant
par une prémisse voisine de l'élimination
déplacée. Sous ces conditions, la diminution
des segments visés n'est compensée ni par
une duplication ni par l'allongement d'un
autre segment de rang maximal.
La seule diminution locale dessinée ne
suffirait pas à établir cette conclusion
pour une permutation arbitraire.













\end{document}
